\chapter{Tungsten and Code Generation}
\label{ch:tungsten}

\begin{chapterintro}
Below the plan, a stage does not interpret operators row by row. Tungsten stores
rows as packed off-heap bytes, and whole-stage code generation collapses an
operator subtree into a single generated Java method that the JIT compiles to
tight machine code. This chapter opens that layer: the binary row format,
off-heap memory management, cache-aware computation, and the generated code you
can actually read with \texttt{EXPLAIN CODEGEN}.
\end{chapterintro}

\section{The Tungsten binary row format}
\tbd

\section{Off-heap memory management}
\tbd

\section{Cache-aware computation}
\tbd

\section{Whole-stage code generation internals}
\tbd

\section{Expression code generation}
\tbd

\section{Columnar processing and vectorization}
\tbd

\section{Summary}
\tbd

\exercisehead
\begin{exercises}
\item Run \texttt{EXPLAIN CODEGEN} on a filter-and-project query. Identify which
      operators were fused into one generated method and which were not.
\end{exercises}
